%%%%%%%%%%%%%%%%%%%%%%%%%%%%%%%%%%%%%%%%%
% Beamer Presentation -- ICPSR Causal Inference
% Day 6 (June 22, 2026): Covariance adjustment -- precision without bias
% Source: separated from the June 18 deck (Days 4-5), expanded to a full session
%%%%%%%%%%%%%%%%%%%%%%%%%%%%%%%%%%%%%%%%%

%------------------------------------------------------------------------------------------------
% PACKAGES AND THEMES
%------------------------------------------------------------------------------------------------

\documentclass[table, xcolor = {dvipsnames}, 9pt]{beamer}
\usepackage{tikz}
\usetikzlibrary{calc}
\usetikzlibrary{positioning}
\usetikzlibrary{arrows.meta}
\usetikzlibrary{external}
\mode<presentation> {
\usetheme{metropolis}
\usecolortheme{seagull}
\usefonttheme{professionalfonts}
}

\usepackage{graphicx} % Allows including images
\usepackage{booktabs} % Allows the use of \toprule, \midrule and \bottomrule in tables
\usepackage{float}
\usepackage{tikz}
\usepackage{multirow}
\usepackage{natbib}
\usepackage{hyperref}
\usepackage{diagbox}
\usepackage{makecell}
\usepackage{xparse}
\usepackage{subfig}
\usepackage{amsmath}
\usepackage{amsfonts,amsthm,amsmath,amssymb}
\usepackage{bbm}
\usepackage{bm}
\usepackage{empheq}
\usepackage{pgfplots}
\usepackage{animate}
\usepgfplotslibrary{colorbrewer}
\pgfplotsset{compat=1.17}

\newcommand\mybox[2][]{\tikz[overlay]\node[fill=lightgray,inner sep=2pt, anchor=text, rectangle, rounded corners=1mm,#1] {#2};\phantom{#2}}
\hypersetup{unicode=true,
            bookmarksnumbered=true,
            bookmarksopen=true,
            bookmarksopenlevel=2,
            breaklinks=false,
            pdfborder={0 0 1},
            hypertexnames=false,
            pdfstartview={XYZ null null 1}}
\usepackage{xcolor}
% R code listing style (matches the course's R-code slides)
\usepackage{listings}
\definecolor{codegreen}{rgb}{0,0.6,0}
\definecolor{codegray}{rgb}{0.5,0.5,0.5}
\definecolor{codepurple}{rgb}{0.58,0,0.82}
\definecolor{backcolour}{rgb}{0.95,0.95,0.92}
\lstdefinestyle{Rstyle}{
    language=R,
    backgroundcolor=\color{backcolour},
    commentstyle=\color{codegreen},
    keywordstyle=\color{magenta},
    stringstyle=\color{codepurple},
    basicstyle=\ttfamily\footnotesize,
    breakatwhitespace=false,
    breaklines=true,
    keepspaces=true,
    showspaces=false,
    showstringspaces=false,
    showtabs=false,
    tabsize=2
}
\newcommand\myheading[1]{%
  \par\bigskip
  {\Large\bfseries#1}\par\smallskip}
\newcommand\given[1][]{\:#1\vert\:}
\theoremstyle{plain}
\newtheorem{thm}{Theorem}
\newtheorem{prop}{Proposition\thisthmnumber}
\newtheorem{lem}{Lemma\thisthmnumber}
\newtheorem{cor}{Corollary}
\newtheorem{defin}{Definition}
\newtheorem{algo}{Algorithm}
\newcommand*\diff{\mathop{}\!\mathrm{d}}
\newcommand*\Diff[1]{\mathop{}\!\mathrm{d^#1}}
\newcommand{\bh}[1]{{\color{blue}{#1}}}
\newcommand{\mh}[1]{{\color{magenta}{#1}}}
\newcommand{\thisthmnumber}{}
\newcommand{\tikzmark}[1]{\tikz[baseline,remember picture] \coordinate (#1) {};}
\newcommand*{\QEDA}{\hfill\ensuremath{\blacksquare}}%
\newcommand*{\QEDB}{\hfill\ensuremath{\square}}%
\newcommand{\abs}[1]{\left\lvert #1 \right\rvert}
\DeclareMathOperator{\E}{\rm{E}}
\DeclareMathOperator{\R}{\mathbb{R}}
\DeclareMathOperator{\N}{\mathbb{N}}
\DeclareMathOperator{\Var}{\rm{Var}}
\DeclareMathOperator{\Cov}{\rm{Cov}}
\DeclareMathOperator{\Supp}{\rm{Supp}}
\DeclareMathOperator{\e}{\rm{e}}
\DeclareMathOperator{\F}{\mathcal{F}}
\DeclareMathOperator{\Z}{\mathcal{Z}}
\DeclareMathOperator{\logit}{\rm{logit}}
\DeclareMathOperator{\indep}{{\perp\!\!\!\perp}}
\DeclareMathOperator{\rank}{rank}
\DeclareMathOperator*{\argmin}{arg\,min}
\DeclareMathOperator*{\argmax}{arg\,max}

%------------------------------------------------------------------------
% TITLE PAGE -- Day 6 (June 22, 2026)
%-----------------------------------------------------------------------
\pagestyle{empty}
\title[Covariance adjustment]{Covariance Adjustment in Randomized Experiments}

\author{Thomas Leavitt}
\institute[]
{
\medskip
\textit{}
}
\date{June 22, 2026}

\begin{document}

\begin{frame}
\titlepage
\end{frame}

%------------------------------------------------------------------------
\section{Where we are}

\begin{frame}[t]{Days 4--5: unbiased estimator, honest inference}
\vfill
\begin{itemize} \vfill
\item \textbf{Estimand}: the ATE $\bar{\tau} = \bar{y}(1) - \bar{y}(0)$ \vfill
\item \textbf{Estimator}: Difference-in-Means $\hat{\tau}\left(\bm{Z}, \bm{Y}\right)$, treated mean $-$ control mean \vfill
\item Under CRA, \emph{whatever} the potential outcomes, $\E\left[\hat{\tau}\left(\bm{Z}, \bm{Y}\right)\right] = \bar{\tau}$ \vfill
\item \mh{Conservative} variance $\Rightarrow$ standard error, test, confidence interval \vfill
\item ACORN: $\hat{\tau} \approx 0.036$, $\widehat{\text{se}} \approx 0.024$, $95\%$ CI includes $0$ \vfill
\end{itemize} \vfill
\end{frame}
%------------------------------------------------------------------------
\begin{frame}[t]{Today: precision from baseline covariates}
\vfill
\begin{itemize} \vfill
\item The Difference-in-Means is unbiased --- but it can be \mh{noisy} \vfill
\item A \mh{baseline covariate} that predicts the outcome carries information \vfill
\item Today: use it to \mh{cut the variance}, leaving unbiasedness intact \vfill
\item Two design-based levers: \vfill
\begin{itemize} \vfill
\item \mh{block} on the covariate before randomizing \vfill
\item \mh{rescale} the outcome by the covariate, including regression adjustment \vfill
\end{itemize} \vfill
\item Running example throughout: \mh{ACORN} GOTV \citep{arceneaux2005} \vfill
\end{itemize} \vfill
\end{frame}
%------------------------------------------------------------------------
\section{Why adjust}

\begin{frame}[t]{Covariance adjustment: precision, not bias}
\vfill
\begin{itemize} \vfill
\item Difference-in-Means is \emph{already} unbiased under CRA --- nothing to fix there \vfill
\item But it can be \mh{noisy} when covariates predict the outcome \vfill
\item \mh{Covariance adjustment}: use baseline covariates to cut the variance \vfill
\end{itemize} \vfill
\end{frame}
%------------------------------------------------------------------------
\begin{frame}[t]{Baseline covariates carry outcome information}
\vfill
\begin{itemize} \vfill
\item \textbf{Baseline covariate}: pre-treatment, so \mh{fixed} regardless of assignment \vfill
\item ACORN: $2002$ turnout predicts potential $2003$ turnout \vfill
\item $x_i$ provides signal about $y_i$ that the raw Difference-in-Means leaves unused \vfill
\end{itemize} \vfill
\end{frame}
%------------------------------------------------------------------------
\begin{frame}[t]{Two design-based ways to use the covariate}
\vfill
\begin{itemize} \vfill
\item \mh{Block}: randomize \emph{within} blocks of similar precincts \vfill
\item \mh{Rescale}: subtract the covariate before testing (a gain score) \vfill
\item Both keep inference \mh{design-based} --- no outcome model is assumed correct \vfill
\item[] each only \emph{tightens} the randomization distribution of $\hat{\tau}$ \vfill
\end{itemize} \vfill
\end{frame}
%------------------------------------------------------------------------
\section{Blocking}

\begin{frame}[t]{If ACORN had been a blocked experiment}
\vfill
\begin{itemize} \vfill
\item Pair the $28$ precincts on $2002$ turnout, $v_{\text{g}2002}$ --- $14$ blocks of two \vfill
\item Randomize \emph{within} each pair: exactly one precinct canvassed \vfill
\item CRA: $\binom{28}{14} \approx 40{,}116{,}600$ assignments \vfill
\item Blocked: $2^{14} = 16{,}384$ assignments --- a tiny, $v_{\text{g}2002}$-balanced subset \vfill
\item \emph{Ideally}, it drops imbalanced assignments --- estimates far from $\bar{\tau}$ \vfill
\item[] So the retained estimates sit nearer the truth $\Rightarrow$ tighter distribution \vfill
\end{itemize} \vfill
\end{frame}
%------------------------------------------------------------------------
\begin{frame}[fragile]{Forming the pairs in \texttt{R} (\texttt{blockTools})}
\begin{lstlisting}[style=Rstyle, basicstyle=\ttfamily\tiny, frame=single]
library(blockTools)
acorn$unit <- seq_len(nrow(acorn))

pairs_out  <- block(
  data       = acorn,        # the data frame of units
  id.vars    = "unit",       # column identifying each unit
  block.vars = "v_g2002",    # covariate(s) to match on
  n.tr       = 2,            # units per block (2 -> matched pairs)
  algorithm  = "optimal")    # minimize total within-pair distance
pair_table <- pairs_out$blocks[[1]]   # 14 rows: "Unit 1", "Unit 2"

# Block randomization: one treated per pair
z_blocked <- integer(nrow(acorn))
for (b in seq_len(nrow(pair_table))) {
  pair <- as.integer(pair_table[b, c("Unit 1", "Unit 2")])
  z_blocked[sample(x = pair, size = 1)] <- 1
}
\end{lstlisting}
\centering \footnotesize \emph{Illustration only}: ACORN was completely randomized; we re-randomize \emph{as if} blocked to show the variance shrink \normalsize
\end{frame}
%------------------------------------------------------------------------
\begin{frame}{Blocking tightens the null distribution}
\begin{figure}[H]
\includegraphics[width=0.92\linewidth]{figures/acorn_blocked_vs_complete.pdf}
\end{figure}
\vspace{-0.3em}
\centering \footnotesize Null SD $0.025 \to 0.021$: the blocked distribution is tighter, hence smaller se \normalsize
\end{frame}
%------------------------------------------------------------------------
\section{Rescaling: the gain score}

\begin{frame}[t]{A fixed covariate shift keeps the estimator unbiased}
\small
\vfill
\begin{itemize} \vfill
\item Gain score $e_i = y_i - x_i$, with POs $e_i(1) = y_i(1) - x_i$, $e_i(0) = y_i(0) - x_i$ \vfill
\item Apply the Difference-in-Means to $\bm{e}$; by Day~4's unbiasedness argument:
\begin{align*}
\E\left[\hat{\tau}\left(\bm{Z}, \bm{e}\right)\right]
 &= \bar{e}(1) - \bar{e}(0)
 && \text{\footnotesize (unbiased for $\bm{e}$'s ATE)} \\
 &= \frac{1}{N}\sum_i \left(y_i(1) - x_i\right) - \frac{1}{N}\sum_i \left(y_i(0) - x_i\right)
 && \text{\footnotesize (def.\ of $e_i$)} \\
 &= \frac{1}{N}\sum_i y_i(1) - \frac{1}{N}\sum_i y_i(0) = \mh{\bar{\tau}}
 && \text{\footnotesize ($x_i$ cancel)}
\end{align*} \vfill
\item Same for any \emph{fixed} $\beta$ in $y_i - \beta x_i$: a pre-treatment $\beta x_i$ cancels in each $\tau_i$ \vfill
\item Needs a \mh{fixed} shift --- estimating $\beta$ from the data breaks \mh{unbiasedness} (next) \vfill
\end{itemize} \vfill
\end{frame}
%------------------------------------------------------------------------
\begin{frame}[fragile]{ACORN in \texttt{R}: covariance adjustment}
\begin{lstlisting}[style=Rstyle, basicstyle=\ttfamily\tiny, frame=single]
x_cov <- acorn$v_g2002              # 2002 general turnout (pre-treatment)
gain_score <- y_obs - x_cov         # rescaled outcome

tau_hat_adj <- diff_in_means(z = z_obs, y = gain_score)            # ~ 0.057
se_hat_adj  <- sqrt(conservative_var(z = z_obs, y = gain_score))   # ~ 0.020
\end{lstlisting}
\begin{itemize} \vfill
\item Standard error $0.024 \to 0.020$: variance about $30\%$ smaller \vfill
\item $95\%$ CI $\to [0.018,\ 0.097]$ --- now \mh{excludes} $0$ (two-sided $p \approx 0.005$) \vfill
\item Estimate moved ($0.036 \to 0.057$): $\hat{\tau}\left(\bm{Z},\bm{e}\right) = \hat{\tau}\left(\bm{Z},\bm{y}\right) - \hat{\tau}\left(\bm{Z},\bm{x}\right)$ \vfill
\end{itemize}
\end{frame}
%------------------------------------------------------------------------
\begin{frame}{Rescaling tightens the null distribution}
\begin{figure}[H]
\includegraphics[width=0.92\linewidth]{figures/acorn_original_vs_rescaled.pdf}
\end{figure}
\vspace{-0.3em}
\centering \footnotesize The gain score's distribution is tighter $\Rightarrow$ smaller se, narrower CI \normalsize
\end{frame}
%------------------------------------------------------------------------
\begin{frame}{The estimate moves, but vanishingly as $N$ grows}
\begin{itemize}
\item The Difference-in-Means is \mh{linear} in the outcome, so $\hat{\tau}\left(\bm{Z}, \bm{y} - \bm{x}\right) = \hat{\tau}\left(\bm{Z}, \bm{y}\right) - \hat{\tau}\left(\bm{Z}, \bm{x}\right)$:
\begin{equation*}
\underbrace{\hat{\tau}\left(\bm{Z}, \bm{e}\right)}_{\text{adjusted}} - \underbrace{\hat{\tau}\left(\bm{Z}, \bm{y}\right)}_{\text{unadjusted}} = -\,\hat{\tau}\left(\bm{Z}, \bm{x}\right), \quad \E\!\left[\hat{\tau}\left(\bm{Z}, \bm{x}\right)\right] = 0,\ \Var \to 0
\end{equation*}
\item[]
\begin{figure}[H]
\includegraphics[width=0.62\linewidth]{figures/acorn_adjust_difference.pdf}
\end{figure}
\item[] \footnotesize Both unbiased for $\bar{\tau}$; the gap $-\hat{\tau}(\bm{Z},\bm{x})$ concentrates at $0$ --- they \mh{agree} asymptotically \normalsize
\end{itemize}
\end{frame}
%------------------------------------------------------------------------
\begin{frame}[t]{What adjustment does --- and does not --- do}
\vfill
\begin{itemize} \vfill
\item \textbf{Does}: cut variance $\Rightarrow$ more \mh{powerful} tests and \mh{tighter} confidence intervals \vfill
\item \textbf{Does not}: remove bias --- randomization gives unbiasedness regardless \vfill
\item \textbf{Does not}: assume correct outcome model --- inference stays \mh{design-based} \vfill
\end{itemize} \vfill
\end{frame}
%------------------------------------------------------------------------
\section{Regression adjustment}

\begin{frame}[t]{Beyond the gain score: regression adjustment}
\vfill
\begin{itemize} \vfill
\item Gain score \emph{fixes} $\beta = 1$ in $y_i - \beta x_i$ --- but why exactly $1$? \vfill
\item Instead residualize with a \mh{least-squares} slope $\hat{\beta}$ \vfill
\item $\hat{\beta}_1, \hat{\beta}_0$: slopes from regressing $y$ on $x$ \emph{within} treated, control \citep{lin2013} \vfill
\item Residualize each unit with its own arm's slope, on centered $x$: \vfill
\begin{equation*}
\tilde{e}_i = y_i - \hat{\beta}_{\mh{Z_i}}\left(x_i - \bar{x}\right), \quad \hat{\beta}_{Z_i} = \begin{cases} \hat{\beta}_1 & Z_i = 1 \\ \hat{\beta}_0 & Z_i = 0 \end{cases}, \quad \hat{\tau}\left(\bm{Z}, \tilde{\bm{e}}\right)
\end{equation*} \vfill
\item Equivalently $\hat{\tau}\left(\bm{Z}, \tilde{\bm{e}}\right) = $ coeff.\ on $Z_i$ in $Y_i \sim Z_i + (\bm{x}_i - \bar{\bm{x}}) + Z_i(\bm{x}_i - \bar{\bm{x}})$  \vfill
\end{itemize} \vfill
\end{frame}
%------------------------------------------------------------------------
\begin{frame}[fragile]{Lin's estimator in \texttt{R}, and what it buys}
\begin{lstlisting}[style=Rstyle, basicstyle=\ttfamily\tiny, frame=single]
X <- scale(acorn[, c("v_g2002","v_p2002","v_m2002")], scale = FALSE)  # center covs
# (a) residualize with arm-specific OLS slopes, then Difference-in-Means
b1 <- coef(lm(vote03 ~ X, data = acorn, subset = z == 1))[-1]  # treated slopes
b0 <- coef(lm(vote03 ~ X, data = acorn, subset = z == 0))[-1]  # control slopes
e  <- acorn$vote03 - ifelse(acorn$z == 1, X %*% b1, X %*% b0)
mean(e[acorn$z == 1]) - mean(e[acorn$z == 0])          # 0.0566
# (b) coefficient on z in the interacted regression -- identical
coef(lm(vote03 ~ z * X, data = acorn))["z"]            # 0.0566
\end{lstlisting}
\begin{itemize} \vfill
\item \mh{Precision}: asymptotically no less precise than unadjusted \citep{lin2013} \vfill
\item \mh{Inference}: HC2 robust se $\overset{p}{\to}$ a limit $\geq$ true se $\Rightarrow$ (conservatively) valid \vfill
\item \mh{Small-sample bias}: estimating $\hat{\beta}$ adds a small finite-sample bias \citep{lin2013} \vfill
\begin{itemize} \vfill
\item[$-$] new research derives exact bias corrections \citep{changetal2024} \vfill
\end{itemize} \vfill
\item \mh{Nonlinear / generalized models}: \citet{guobasse2022}, \citet{cohenfogarty2022b} \vfill
\end{itemize}
\end{frame}
%------------------------------------------------------------------------
\section{Recap}

\begin{frame}[t]{Recap: precision by design}
\vfill
\begin{itemize} \vfill
\item A \mh{baseline} covariate is fixed under both potential outcomes \vfill
\item Cut variance, \emph{not} bias --- randomization already $\implies$ unbiasedness \vfill
\item \mh{Block} on it, or \mh{rescale} (gain score), then regression-adjust \citep{lin2013} \vfill
\item ACORN: se $0.024 \to 0.020$; the $95\%$ CI now excludes $0$ \vfill
\item Inference stays \mh{design-based}: no outcome model assumed correct \vfill
\item[] next: when assignment $\neq$ receipt --- \mh{noncompliance} \vfill
\end{itemize} \vfill
\end{frame}
%------------------------------------------------------------------------
\begin{frame}[t]{In-class exercise}
\vfill
\begin{itemize} \vfill
\item With a neighbor, using the ACORN data: \vfill
\begin{enumerate} \vfill
\item Recompute the gain-score estimate and se using \texttt{v\_p2002} instead of \texttt{v\_g2002}. Does precision improve as much? \vfill
\item Run \texttt{lm\_lin(vote03 \textasciitilde\ z, covariates = \textasciitilde\ v\_g2002)}. Compare its estimate and HC2 se to the gain-score Difference-in-Means and its se. \vfill
\item Why does the adjusted estimate move from $0.036$ to $0.057$, even though both the adjusted and unadjusted \emph{estimators} are unbiased for $\bar{\tau}$? \vfill
\end{enumerate} \vfill
\end{itemize} \vfill
\end{frame}
%------------------------------------------------------------------------
\begin{frame}[allowframebreaks]
\frametitle{References}
\scriptsize
\bibliographystyle{chicago}
\bibliography{Master_Bibliography}
\end{frame}
%------------------------------------------------------------------------
\end{document}
